\chapter{Preliminaries}\label{ch:prelim}

This chapter fixes the categorical and algebraic vocabulary used throughout
the thesis. Everything stated here is standard: definitions are given, the
sources of record are cited, and no proofs are included. The presentation is
self-contained in the sense that every notion appealed to in later chapters
is defined here or directly cited.

\section{Categories}\label{sec:prelim-cat}

\begin{definition}[Category]
A \emph{category} $\mathcal{C}$ consists of
\begin{itemize}
  \item a collection $\mathrm{Ob} (\mathcal{C})$ of \emph{objects};
  \item for every pair of objects $A, B$, a set $\mathcal{C} (A,B)$ of
        \emph{morphisms}, the \emph{hom-set}, with $f \in \mathcal{C} (A,B)$
        written $f \colon A \to B$;
  \item for every object $A$, an \emph{identity} morphism
        $1_A \in \mathcal{C} (A,A)$;
  \item for every triple of objects $A, B, C$, a \emph{composition} function
        \[ \circ \colon \mathcal{C} (B,C) \times \mathcal{C} (A,B)
           \to \mathcal{C} (A,C), \qquad (g, f) \mapsto g \circ f, \]
\end{itemize}
such that composition is \emph{associative}, $(h \circ g) \circ f
= h \circ (g \circ f)$, and the \emph{identity laws} hold,
$1_B \circ f = f = f \circ 1_A$~\cite{maclane}.
\end{definition}

\begin{definition}[Isomorphism]
A morphism $f \colon A \to B$ in a category is an \emph{isomorphism} if
there exists a morphism $g \colon B \to A$ with $g \circ f = 1_A$ and
$f \circ g = 1_B$; the objects $A$ and $B$ are then \emph{isomorphic},
written $A \cong B$~\cite{maclane}.
\end{definition}

\begin{remark}
A category is \emph{small} if its objects and morphisms form sets and
\emph{locally small} if every hom-set $\mathcal{C} (A,B)$ is a set; all
categories in this thesis are locally small, and the distinction matters only
for \S\ref{sec:prelim-acc}~\cite{maclane}.
\end{remark}

\begin{example}
The category $\mathbf{Set}$ of sets and total functions is the running
example. Reading the objects of a typed language as types, the total
functions between them form a category; this is the discrete approximation of
the category $\mathbf{PureMol}$ of \autoref{ch:mol}.
\end{example}

\section{Functors}\label{sec:prelim-func}

\begin{definition}[Functor]
A \emph{functor} $F \colon \mathcal{C} \to \mathcal{D}$ assigns to each
object $A$ of $\mathcal{C}$ an object $F A$ of $\mathcal{D}$ and to each
morphism $f \colon A \to B$ a morphism $F f \colon F A \to F B$, preserving
identities and composition:
$F 1_A = 1_{F A}$ and $F (g \circ f) = F g \circ F f$~\cite{maclane}.
\end{definition}

\begin{definition}[Full, faithful, essentially surjective]
A functor $F \colon \mathcal{C} \to \mathcal{D}$ is \emph{full} if
$\mathcal{C} (A,B) \to \mathcal{D} (F A, F B)$ is surjective for all $A, B$,
\emph{faithful} if it is injective for all $A, B$, and \emph{essentially
surjective on objects} if every object of $\mathcal{D}$ is isomorphic to
$F A$ for some object $A$ of $\mathcal{C}$~\cite{maclane}.
\end{definition}

\begin{definition}[Equivalence of categories]
An \emph{equivalence} of categories is a functor that is full, faithful, and
essentially surjective on objects; we write $\mathcal{C} \simeq
\mathcal{D}$~\cite{maclane}.
\end{definition}

\section{Natural Transformations}\label{sec:prelim-nat}

\begin{definition}[Natural transformation]
Let $F, G \colon \mathcal{C} \to \mathcal{D}$ be functors. A \emph{natural
transformation} $\alpha \colon F \Rightarrow G$ is a family of morphisms
$\alpha_A \colon F A \to G A$, one for each object $A$ of $\mathcal{C}$, such
that for every morphism $f \colon A \to B$ the \emph{naturality condition}
\[ G f \circ \alpha_A = \alpha_B \circ F f \]
holds, i.e.\ the square built from $\alpha_A$, $\alpha_B$, $F f$, and $G f$
commutes~\cite{maclane}.
\end{definition}

\begin{definition}[Natural isomorphism]
A natural transformation $\alpha \colon F \Rightarrow G$ is a \emph{natural
isomorphism} if every component $\alpha_A$ is an isomorphism; the functors
are then \emph{naturally isomorphic}, $F \cong G$~\cite{maclane}.
\end{definition}

\begin{remark}
Functors and natural transformations organize (small) categories into a
category $\mathbf{Cat}$. Its isomorphisms are the invertible functors, while
the more useful notion of sameness is the equivalence of categories: a
functor $F$ is an equivalence exactly when there is a functor $G$ with
$G \circ F \cong 1_{\mathcal{C}}$ and $F \circ G \cong 1_{\mathcal{D}}$
(natural isomorphisms). This is the relation asserted between
$\mathbf{PureMol}$ and the category of types and total functions in
\autoref{ch:mol}.
\end{remark}

\section{Monoidal and Symmetric Monoidal Categories}\label{sec:prelim-mon}

\begin{definition}[Monoidal category]
A \emph{monoidal category} is a category $\mathcal{C}$ equipped with
\begin{itemize}
  \item a \emph{tensor product} bifunctor
        $\otimes \colon \mathcal{C} \times \mathcal{C} \to \mathcal{C}$,
        written $(A,B) \mapsto A \otimes B$;
  \item a \emph{unit object} $I$;
  \item natural isomorphisms $\alpha_{A,B,C} \colon (A \otimes B) \otimes C
        \to A \otimes (B \otimes C)$ (the \emph{associator}) and
        $\lambda_A \colon I \otimes A \to A$, $\rho_A \colon A \otimes I \to A$
        (the \emph{left} and \emph{right unitors}),
\end{itemize}
such that the \emph{pentagon} and \emph{triangle} coherence diagrams commute;
e.g.\ the pentagon states the equality of the two morphisms
$((A \otimes B) \otimes C) \otimes D \to A \otimes (B \otimes (C \otimes D))$
built from $\alpha$ and the tensor of $\alpha$ with identities~\cite{maclane}.
\end{definition}

\begin{definition}[Strict monoidal category]
A monoidal category is \emph{strict} if $\alpha$, $\lambda$, and $\rho$ are
identity natural transformations, so that $(A \otimes B) \otimes C
= A \otimes (B \otimes C)$ and $I \otimes A = A = A \otimes I$ on the nose.
\end{definition}

\begin{definition}[Symmetric monoidal category]
A \emph{symmetric monoidal category} is a monoidal category with a natural
isomorphism $\sigma_{A,B} \colon A \otimes B \to B \otimes A$ (the
\emph{symmetry}) such that $\sigma_{B,A} \circ \sigma_{A,B} = 1_{A \otimes B}$
and $\sigma$ is coherent with the associator and the unitors~\cite{maclane}.
\end{definition}

\begin{remark}[Coherence]
Mac Lane's coherence theorem states that every diagram whose edges are built
from $\alpha$, $\lambda$, $\rho$, and $\sigma$ commutes; hence no information
is lost in treating $A \otimes (B \otimes C)$ and $(A \otimes B) \otimes C$
as interchangeable. Consequently every monoidal category is monoidally
equivalent to a strict one (Mac Lane's strictification), a fact used
implicitly throughout~\cite{maclane}. The tensor product of \autoref{ch:mol}
is tuple formation with unit $()$, which is strict and symmetric.
\end{remark}

\section{Traced Monoidal Categories}\label{sec:prelim-trace}

\begin{definition}[Traced monoidal category]
Let $\mathcal{C}$ be a symmetric monoidal category. A \emph{trace} is a
family of operations
\[ \mathrm{Tr}^{U}_{A,B} \colon \mathcal{C} (A \otimes U,\, B \otimes U)
   \to \mathcal{C} (A,B), \]
natural in $A$ and $B$ and dinatural in the \emph{feedback object} $U$,
satisfying the following axioms~\cite{joyalsv, arxiv-traced-monoidal}:
\begin{itemize}
  \item \emph{Vanishing}: $\mathrm{Tr}^{I}_{A,B} (f) = f$ for
        $f \colon A \otimes I \to B \otimes I$, and
        $\mathrm{Tr}^{U \otimes V}_{A,B}
        = \mathrm{Tr}^{U}_{A,B} \circ
          \mathrm{Tr}^{V}_{A \otimes U, B \otimes U}$;
  \item \emph{Superposing}: for $f \colon A \otimes U \to B \otimes U$ and
        $g \colon C \to D$,
        \[ \mathrm{Tr}^{U}_{A \otimes C, B \otimes D} (f \otimes g)
           = \mathrm{Tr}^{U}_{A,B} (f) \otimes g; \]
  \item \emph{Yanking}: $\mathrm{Tr}^{U}_{U,U} (\sigma_{U,U}) = 1_U$.
\end{itemize}
A \emph{traced monoidal category} is a symmetric monoidal category equipped
with a trace. The trace internalizes feedback: a morphism
$f \colon A \otimes U \to B \otimes U$ computes a transformation $A \to B$
while threading the feedback object $U$ through the loop. This is the
structure underlying the fixed-iteration loops of \autoref{ch:trace}.
\end{definition}

\section{Monads, Algebras over Monads, and Kleisli Categories}\label{sec:prelim-monads}

\begin{definition}[Monad]
A \emph{monad} on a category $\mathcal{C}$ is a triple $(T, \eta, \mu)$ where
$T \colon \mathcal{C} \to \mathcal{C}$ is a functor and
$\eta \colon 1_{\mathcal{C}} \Rightarrow T$ (the \emph{unit}) and
$\mu \colon T^2 \Rightarrow T$ (the \emph{multiplication}) are natural
transformations satisfying the \emph{unit laws}
$\mu \circ T\eta = 1_T = \mu \circ \eta T$ and the \emph{associativity law}
$\mu \circ T\mu = \mu \circ \mu T$~\cite{maclane, arxiv-monads-kleisli}.
\end{definition}

\begin{definition}[Algebra over a monad]
Let $(T, \eta, \mu)$ be a monad on $\mathcal{C}$. A \emph{$T$-algebra}
(an \emph{Eilenberg--Moore algebra}) is a pair $(A, a)$ consisting of an
object $A$ and a morphism $a \colon T A \to A$ such that
$a \circ \eta_A = 1_A$ and $a \circ \mu_A = a \circ T a$. A \emph{homomorphism}
of $T$-algebras $(A,a) \to (B,b)$ is a morphism $f \colon A \to B$ with
$f \circ a = b \circ T f$. $T$-algebras and their homomorphisms form the
category $\mathcal{C}^T$~\cite{maclane}.
\end{definition}

\begin{definition}[Kleisli category]
The \emph{Kleisli category} $\mathcal{C}_T$ of a monad $(T, \eta, \mu)$ on
$\mathcal{C}$ has the same objects as $\mathcal{C}$; a morphism $A \to B$ in
$\mathcal{C}_T$ is a morphism $f \colon A \to T B$ in $\mathcal{C}$. The
identity on $A$ is $\eta_A$, and the composition of $f \colon A \to T B$ with
$g \colon B \to T C$ is
\[ g \circ_T f = \mu_C \circ T g \circ f \colon A \to T C, \]
the \emph{bind} operation of the monad~\cite{maclane, arxiv-monads-kleisli}.
\end{definition}

\begin{example}[State monad]
Let $\mathcal{C}$ have finite products and fix an object $S$ of
\emph{states}. The \emph{state monad} is the monad with
$T X = (X \times S)^S$, unit $\eta_X \colon x \mapsto \lambda s.\,(x, s)$,
and multiplication given by threading the state through two successive
computations. A Kleisli morphism $A \to B$ is then a morphism
$A \times S \to B \times S$, exactly the transition function of a
deterministic Mealy machine with state space $S$
~\cite{maclane, arxiv-monads-kleisli, mealy, moore, arxiv-mealy}. The
Kleisli category of this monad models $\mathbf{EffMol} (\mathrm{Ctx})$ in
\autoref{ch:mol}, and the pure/effectful decomposition of
\autoref{ch:kleisli} is phrased in this language.
\end{example}

\begin{remark}
Every monad arises from an adjunction; the Kleisli category is the initial
and $\mathcal{C}^T$ the terminal resolution of the monad as a pair of
adjoint functors~\cite{maclane, arxiv-monads-kleisli}.
\end{remark}

\section{Coalgebras}\label{sec:prelim-coalg}

\begin{definition}[Coalgebra for an endofunctor]
Let $F \colon \mathcal{C} \to \mathcal{C}$ be an endofunctor. An
\emph{$F$-coalgebra} is a pair $(S, s)$ consisting of an object $S$ and a
morphism $s \colon S \to F S$. A \emph{homomorphism} of $F$-coalgebras
$(S,s) \to (T,t)$ is a morphism $f \colon S \to T$ with
$t \circ f = F f \circ s$. Coalgebras and their homomorphisms form the
category $\mathbf{Coalg} (F)$~\cite{arxiv-coalgebra-min}.
\end{definition}

\begin{definition}[Final coalgebra]
A \emph{final} (terminal) coalgebra is a terminal object of
$\mathbf{Coalg} (F)$: a coalgebra $(Z, \zeta)$ such that for every coalgebra
$(S,s)$ there is a unique homomorphism $(S,s) \to (Z,\zeta)$~\cite{arxiv-coalgebra-min}.
\end{definition}

\begin{definition}[Behavioral equivalence]
Two states $s_1 \in S_1$, $s_2 \in S_2$ of $F$-coalgebras are
\emph{behaviorally equivalent} if there is a coalgebra $(T,t)$ and
homomorphisms from $(S_1,s_1)$ and $(S_2,s_2)$ into it identifying them; when
$F$ preserves weak pullbacks this coincides with \emph{bisimilarity} in the
sense of~\cite{park, arxiv-coalgebra-min}.
\end{definition}

\begin{example}[Mealy machines as coalgebras]
Fix sets $A$ and $B$ of inputs and outputs. A deterministic Mealy machine
$(S, \mathrm{step})$ with $\mathrm{step} \colon S \times A \to B \times S$
is exactly a coalgebra for the functor $F X = (B \times X)^A$
~\cite{mealy, moore, arxiv-mealy}. Behavioral equivalence is then language
equivalence, and the quotient constructions of \autoref{ch:coalgebra}
refine the classical Nerode and Hopcroft minimization
algorithms~\cite{nerode, hopcroft, arxiv-coalgebra-min}.
\end{example}

\section{Lawvere Theories}\label{sec:prelim-lawvere}

\begin{definition}[Lawvere theory]
Let $\mathbf{Fin}^{\mathrm{op}}$ be the category whose objects are the
natural numbers and whose morphisms $m \to n$ are the functions $n \to m$.
A \emph{Lawvere theory} is a category $\mathcal{L}$ with finite products
together with an identity-on-objects, strict, finite-product-preserving
functor $\mathbf{Fin}^{\mathrm{op}} \to \mathcal{L}$. The object $1$ is the
\emph{generic object}: every object of $\mathcal{L}$ is a finite product of
copies of $1$~\cite{lawvere, arxiv-lawvere}.
\end{definition}

\begin{definition}[Model of a Lawvere theory]
A \emph{model} of a Lawvere theory $\mathcal{L}$ in a category $\mathcal{C}$
with finite products is a finite-product-preserving functor
$\mathcal{L} \to \mathcal{C}$. Models for $\mathcal{C} = \mathbf{Set}$ are
the \emph{algebras} of the theory, and models and natural transformations
form a category $\mathbf{Mod} (\mathcal{L}, \mathcal{C})$~\cite{lawvere,
arxiv-lawvere}.
\end{definition}

\begin{remark}
The \emph{free algebra} on $n$ generators is the model represented by object
$n$: the finite-product-preserving functor $\mathcal{L} (n,-)$ with values in
$\mathbf{Set}$, whose underlying set is the set of terms in $n$ variables
modulo the equations of the theory~\cite{lawvere}. A finitary equational
theory in the sense of universal algebra (Section~\ref{sec:prelim-univ}) is
the same data as a Lawvere theory, and its variety of algebras is
$\mathbf{Mod} (\mathcal{L}, \mathbf{Set})$. The equations of the dataplane
component theories of \autoref{ch:lawvere} (ring, parser, transport) are
Lawvere theories whose models are the corresponding molecules.
\end{remark}

\section{Operads, Multicategories, and Clubs}\label{sec:prelim-operads}

\begin{definition}[Symmetric operad]
A \emph{symmetric operad} $P$ (in $\mathbf{Set}$) consists of a set $P (n)$
of \emph{$n$-ary operations} for each $n \geq 0$, an action of the symmetric
group $\Sigma_n$ on $P (n)$, a distinguished element $1 \in P (1)$ (the
\emph{identity operation}), and composition maps
\[ P (k) \times P (n_1) \times \cdots \times P (n_k)
   \to P (n_1 + \cdots + n_k), \]
that are associative, unital with respect to $1$, and equivariant with
respect to the symmetric group actions~\cite{leinster, arxiv-operads}.
\end{definition}

\begin{definition}[Algebra over an operad]
An \emph{algebra} over a symmetric operad $P$ is a set $A$ together with, for
each $p \in P (n)$, an operation $A^n \to A$, respecting the identity, the
composition maps, and the symmetric group action~\cite{leinster,
arxiv-operads}.
\end{definition}

\begin{definition}[Multicategory]
A \emph{multicategory} is a category-like structure in which a morphism has a
finite \emph{sequence} of objects as domain: a morphism
$(A_1, \dots, A_n) \to B$, with identity morphisms $1_A \colon (A) \to A$ and
an associative \emph{substitution} composition that is unital. A
\emph{symmetric} multicategory carries a compatible action of $\Sigma_n$ on
morphisms with $n$-ary domain~\cite{leinster, arxiv-operads}.
\end{definition}

\begin{remark}
An operad is precisely a one-object symmetric multicategory: its morphisms
whose domain is a list of copies of the single object are the operations, and
substitution is operadic composition. Kelly's \emph{clubs} generalize both: a
club is a category over the category of finite sets (or of permutations)
whose substitution structure presents the free algebras of a 2-monad,
covering many doctrines of categories with structure, including the symmetric
monoidal ones at work in \autoref{ch:operad} and
\autoref{ch:rewrite}~\cite{kelly, leinster}.
\end{remark}

\section{Effectus Theory}\label{sec:prelim-effectus}

Effectus theory formalizes partiality, probabilistic choice, and measurement
in one categorical framework, axiomatically oriented around the
\emph{effects} (observables) of a system.

\begin{definition}[Effect algebra]
An \emph{effect algebra} is a set $E$ with a partial binary operation
$\oplus$ and distinguished elements $0 \neq 1$, such that $\oplus$ is
commutative and associative wherever defined, $a \oplus 0 = a$ for all $a$,
every element $a$ has a unique \emph{orthosupplement} $a^{\perp}$ with
$a \oplus a^{\perp} = 1$, and $a \oplus 1$ is defined only for
$a = 0$~\cite{arxiv-effectus}.
\end{definition}

\begin{definition}[Effectus]
An \emph{effectus} is a category with finite coproducts and a terminal object
$1$, admitting the pushouts of coproduct injections that are needed to
compose \emph{states} and \emph{effects}, such that for every object $X$ the
hom-set of \emph{effects} $X \to 1$ carries an effect algebra structure and
composition induces effect-algebra homomorphisms
~\cite{arxiv-effectus}.
\end{definition}

\begin{remark}
Effects generalize predicates: on a measurable space the effects are the
$[0,1]$-valued measurable functions with $\oplus$ truncated addition. The
framework supplies one axiomatics for partiality and probabilistic choice,
relevant to the error and context channels of the effectful molecules of
\autoref{ch:mol}; a closely related categorical treatment of probability is
developed in the theory of Markov categories
~\cite{arxiv-effectus, arxiv-markov-categories}.
\end{remark}

\section{Linear Logic and Geometry of Interaction}\label{sec:prelim-linear}

\begin{definition}[Linear logic]
\emph{Linear logic} is a substructural logic in which every hypothesis is
used exactly once. It is presented by sequent rules for the multiplicative
connectives $\otimes$ (tensor) and $\multimap$ (linear implication), the
additive connectives $\&$ and $\oplus$, and the \emph{exponentials} $!$ and
$?$ that govern contraction and weakening, i.e.\ unlimited use of a
hypothesis~\cite{girard}.
\end{definition}

\begin{remark}[Cut elimination and geometry of interaction]
The \emph{cut rule} composes proofs, and \emph{cut elimination} rewrites a
proof to a cut-free normal form, which is deterministic and is the
computational content of the proof. The \emph{geometry of interaction} (GoI)
interprets proofs as morphisms of a traced monoidal category and cut
elimination as the trace: the \emph{Int} construction embeds every traced
monoidal category into a compact closed one, where the trace becomes genuine
composition~\cite{girard, joyalsv, arxiv-linear-logic-goi}. Deforestation
and the zero-allocation fusion of \autoref{ch:linear} are the thesis's
instance of this correspondence.
\end{remark}

\section{Accessible and Locally Presentable Categories}\label{sec:prelim-acc}

\begin{definition}[Presentable objects, accessible categories]
Let $\kappa$ be a regular cardinal. An object $A$ of a category
$\mathcal{C}$ is \emph{$\kappa$-presentable} if the hom-functor
$\mathcal{C} (A, -)$ preserves $\kappa$-filtered colimits. The category
$\mathcal{C}$ is \emph{$\kappa$-accessible} if it has $\kappa$-filtered
colimits and there is a set of $\kappa$-presentable objects such that every
object of $\mathcal{C}$ is a $\kappa$-filtered colimit of objects from that
set; $\mathcal{C}$ is \emph{accessible} if it is $\kappa$-accessible for some
regular cardinal $\kappa$~\cite{adamek}.
\end{definition}

\begin{definition}[Locally presentable categories]
A category is \emph{locally $\kappa$-presentable} if it is
$\kappa$-accessible and cocomplete, and \emph{locally presentable} if it is
locally $\kappa$-presentable for some regular cardinal $\kappa$~\cite{adamek}.
\end{definition}

\begin{remark}
Every locally presentable category is complete and cocomplete, and the
categories of models of Lawvere theories (Section~\ref{sec:prelim-lawvere})
and the categories of algebras of many signatures are locally presentable.
Local presentability is the standard tool for proving the existence of free
structures, such as the free symmetric monoidal category $F (\Sigma)$ whose
existence \autoref{ch:free} invokes~\cite{adamek, lawvere}.
\end{remark}

\section{Universal Algebra and Clones}\label{sec:prelim-univ}

\begin{definition}[Signature and algebra]
A \emph{signature} is a set $\Sigma$ of operation symbols together with an
assignment of a finite \emph{arity} to each symbol. A \emph{$\Sigma$-algebra}
is a set $A$ together with, for each $n$-ary symbol $\sigma \in \Sigma$, an
operation $\sigma_A \colon A^n \to A$~\cite{knuthbendix}.
\end{definition}

\begin{definition}[Terms, equations, varieties]
Given a signature $\Sigma$ and a set $X$ of variables, the \emph{terms} over
$X$ are formed from variables and operation symbols. An \emph{equation} is a
pair of terms $s \approx t$; a $\Sigma$-algebra \emph{satisfies} the equation
if the two terms denote equal operations under every assignment of values to
variables. A class of algebras closed under satisfaction of a set of
equations is a \emph{variety}~\cite{lawvere, knuthbendix}.
\end{definition}

\begin{definition}[Clone]
A \emph{clone} on a set $A$ is a set $C$ of finitary operations on $A$ that
contains all projections and is closed under composition. The clone of a
$\Sigma$-algebra is the set of its \emph{term operations}~\cite{lawvere}.
\end{definition}

\begin{remark}
A clone on $A$ is the concrete shadow of a Lawvere theory: it is exactly the
data of a model of an algebraic theory with carrier $A$
(Section~\ref{sec:prelim-lawvere}). The \emph{word problem} for a variety
asks whether two terms denote the same operation in the free algebra; the
Knuth--Bendix completion procedure decides it whenever it terminates, by
producing a terminating and confluent rewriting system, and Newman's lemma
reduces confluence to local confluence~\cite{knuthbendix, newman,
arxiv-rewriting-kb}. By the classical theorem of Birkhoff \citep{birkhoff}, the varieties are
precisely the classes of algebras closed under homomorphic images,
subalgebras, and products. The equational theory of the runtime is treated by
these means in \autoref{ch:rewrite}.
\end{remark}

\begin{table}[htbp]
\centering
\small
\begin{tabular}{@{}lll@{}}
\toprule
Notion & Defined in & Applied in \\
\midrule
Categories, functors, natural transformations & \S\ref{sec:prelim-cat} & \autoref{ch:mol} \\
Monoidal and symmetric monoidal categories & \S\ref{sec:prelim-mon} & \autoref{ch:mol}, \autoref{ch:free} \\
Traced monoidal categories & \S\ref{sec:prelim-trace} & \autoref{ch:trace} \\
Monads, algebras, Kleisli categories & \S\ref{sec:prelim-monads} & \autoref{ch:kleisli}, \autoref{ch:mol} \\
Coalgebras & \S\ref{sec:prelim-coalg} & \autoref{ch:coalgebra} \\
Lawvere theories & \S\ref{sec:prelim-lawvere} & \autoref{ch:lawvere} \\
Operads, multicategories, clubs & \S\ref{sec:prelim-operads} & \autoref{ch:operad}, \autoref{ch:rewrite} \\
Effectus theory & \S\ref{sec:prelim-effectus} & \autoref{ch:mol} \\
Linear logic and geometry of interaction & \S\ref{sec:prelim-linear} & \autoref{ch:linear} \\
Accessible and locally presentable categories & \S\ref{sec:prelim-acc} & \autoref{ch:free} \\
Universal algebra and clones & \S\ref{sec:prelim-univ} & \autoref{ch:rewrite} \\
\bottomrule
\end{tabular}
\caption{The categorical tools of this chapter and where the thesis applies them.}
\end{table}


\section{Notation and acronyms}\label{sec:prelim-notation}

All acronyms are defined here and at their first use.

\paragraph{Acronyms.}
\begin{itemize}
\item \textbf{SMC}: symmetric monoidal category (\autoref{sec:prelim-mon}).
\item \textbf{DFA}: deterministic finite automaton.
\item \textbf{EMA}: exponentially weighted moving average (the
  scheduler estimate of \autoref{sec:scheduler}).
\item \textbf{IPC}: instructions per cycle, a processor counter.
\item \textbf{TLB}: translation lookaside buffer; a dTLB is the data
  TLB.
\item \textbf{RSS}: resident set size, the physical memory a process
  occupies.
\item \textbf{L1, L3}: first- and last-level processor caches.
\item \textbf{TCP, UDP, SCTP}: the Transmission Control Protocol, the
  User Datagram Protocol, and the Stream Control Transmission Protocol.
\item \textbf{HTTP}: the Hypertext Transfer Protocol.
\item \textbf{NIC}: network interface controller.
\item \textbf{DNS}: the Domain Name System.
\item \textbf{FTP}: the File Transfer Protocol.
\item \textbf{LL(1)}: left-to-right, leftmost-derivation parsing with
  one token of lookahead.
\item \textbf{RFC}: Request for Comments, the standards document series
  of the Internet Engineering Task Force.
\item \textbf{SQPOLL}: submission-queue polling, an io\_uring mode in
  which the kernel polls the submission queue and the application skips
  the enter-system-call.
\item \textbf{TLS}: Transport Layer Security.
\item \textbf{QUIC}: the connection-oriented transport protocol used
  by HTTP/3.
\item \textbf{SIMD}: single instruction, multiple data.
\item \textbf{KB}: Knuth--Bendix completion (\autoref{ch:rewrite}).
\item \textbf{TTL}: time to live, a cache validity window.
\item \textbf{AF\_XDP}: the kernel's eXpress Data Path socket family.
\item \textbf{io\_uring}: the kernel's asynchronous completion-ring
  interface.
\item \textbf{MSG\_ZEROCOPY}: the socket flag for zero-copy sends.
\item \textbf{PGO}: profile-guided optimization.
\end{itemize}

\paragraph{Notation.}
$\circ$ is sequential composition and $\otimes$ is parallel
composition (tensor); $\mathrm{id}_A$ is the identity on the object
$A$; $\simeq$ is behavioral equivalence; $\beta(M)$ is the boundary
count of Definition~\ref{def:boundary}; $c(M)$ is the cost vector of
Definition~\ref{def:cost}; $\mathrm{NF}(M)$ is the normal form of $M$.
